\chapter*{General Conclusion}
\addcontentsline{toc}{chapter}{General Conclusion}

This end-of-studies internship at Adria Business \& Technology provided a direct view of how a banking software platform is designed, implemented, reviewed, and refined inside an enterprise team. The work was not limited to writing isolated code. It was carried by a complete delivery rhythm: Jira tickets, Bitbucket branches, pull requests, technical lead review, tester validation, Scrum ceremonies, and Product Owner refinement sessions.

The report first presented the host organization and the Supply Chain Finance context. This was necessary because the technical platform only makes sense when connected to its business purpose: financing invoices, managing counterparties, configuring programs, enforcing maker/checker governance, and giving banks better operational visibility. The following chapters then examined the system architecture, the backend domain model, the frontend flows, and the Disbursement Module assigned as the academic project.

The Disbursement Module represents the main contribution of the PFE subject. Its design clarified the boundary between invoice lodging, financing eligibility, payment execution, finance-record creation, and repayment. The module also showed how one business operation can cross several technical concerns: entity modeling, frontend wizard design, program limits, asynchronous callbacks, Kafka events, Payment Gateway integration, and transaction inquiry. The work remained partly target-oriented because the platform had not yet reached the sprint planned for native integration of the module. This limitation is important, but it does not erase the value of the analysis. It shows the real state of the project and the preparation needed before implementation becomes final.

The internship also produced a broader learning outcome. It improved technical practice in Spring Boot, React, TypeScript, JPA, OAuth2 security, gateway routing, and service-oriented organization. It also developed professional habits: reading existing code before changing it, accepting review comments, explaining decisions, translating business rules into diagrams, and handling ambiguity without inventing unsupported assumptions.

Future work should continue with the native integration of the Disbursement Module, final validation of the backend contract, stronger end-to-end testing, richer reporting, and more complete asynchronous processing around payments and notifications. The project therefore closes with a working understanding of the platform and a clear continuation path. It is both a technical experience and a professional one.

